\chapter{この教材の歩き方}

\section{ゴール}

この教材のゴールは、\textbf{理工体剣道部の公式サイトを、自分で保守・運用できるようになること}です。
保守というのは、たとえば次のようなことを、三角に頼らず自分で片づけられる状態を指します。

\begin{itemize}
  \item 新しい年度の部員名簿や大会結果を、管理画面から反映できる
  \item 「ここの文言を直して」「写真を差し替えて」と頼まれたら、コードを直して公開できる
  \item 「サイトが表示されない」と言われたとき、どこが壊れているのかを絞り込める
  \item 字幕アプリの処理が失敗したとき、やり直してよいかどうかを判断できる
\end{itemize}

そのために、部のサイトの縮小版を、自分の手で一から作り直します。
完成品はすでにあるので、作り直すことに意味があるのか不思議に思うかもしれません。
けれども、\textbf{自分で一度組み立てたものは、壊れたときに直せる}ようになります。
これがこの教材の一番の狙いです。

\section{全体の地図}

この教材は全6巻あり、いま読んでいるのは第1巻です。

\begin{center}
\begin{tikzpicture}[
  node distance=4mm,
  ph/.style={draw=kon,rounded corners=3pt,fill=white,minimum height=12mm,text width=2.25cm,
    align=center,font=\sffamily\small},
  now/.style={ph,fill=kon,text=white},
  arr/.style={-{Stealth[length=2.5mm]},thick,kon!70}]
  \node[now] (v1) {第1巻\\環境構築\\HTML/CSS/JS};
  \node[ph,right=of v1] (v2) {第2巻\\Next.js で\\作り直す};
  \node[ph,right=of v2] (v3) {第3巻\\DB と公開};
  \node[ph,right=of v3] (v4) {第4巻\\ログインと\\管理画面};
  \node[ph,below=8mm of v4] (v5) {第5巻\\障害対応と\\本番の保守};
  \node[ph,left=of v5] (v6) {第6巻\\字幕アプリ\\の保守};
  \draw[arr] (v1)--(v2); \draw[arr] (v2)--(v3); \draw[arr] (v3)--(v4);
  \draw[arr] (v4)--(v5); \draw[arr] (v5)--(v6);
  \node[font=\sffamily\scriptsize,text=shu,below=1mm of v1] {いまここ};
\end{tikzpicture}
\end{center}

第1巻では、開発に使う道具をそろえたあと、HTML・CSS・JavaScript だけで部のトップページを作ります。
完成したら、インターネットに公開するところまで進めます。

\begin{why}[なぜいきなり Next.js から始めないのか]
本番のサイトは Next.js というフレームワークで作られています。
ただ、Next.js は「HTML・CSS・JavaScript を楽に書くための仕組み」なので、
土台の3つを知らないまま使うと、エラーが出たときに何が起きているのか分かりません。
第1巻で素の HTML・CSS・JavaScript を一度書いておくと、第2巻で Next.js を使ったときに
「あの面倒な作業を、これが代わりにやってくれているのか」と納得しながら進められます。
\end{why}

\section{教材と練習用リポジトリの場所}

この教材と、第2巻から使う練習用のリポジトリは、GitHub の次の場所にあります。

\begin{center}
\fcolorbox{kon}{kon!4}{\parbox{0.85\linewidth}{\centering\sffamily
\url{https://github.com/yutomi141/rikotaikendo-practice}\\[2pt]
{\small 教材の PDF：\code{docs/textbook/vol1.pdf}（第2巻以降も同じ場所に増えていく）}}}
\end{center}

\begin{itemize}
  \item 第1巻のうちは、この PDF を読むだけで大丈夫です。リポジトリへの招待は第2章、手元へのコピー（clone）は第6章で行います。
  \item 教材が新しくなったら、三角が LINE で知らせます。最新の PDF は、いつでも上の場所から取れます。
  \item 教材の元のファイル（\code{.tex}）も同じ場所にあるので、第6章以降は Claude Code に「教材の第7章のここが分からない」と聞くこともできます。
\end{itemize}

\section{この教材の読み方：3つのルール}

\begin{enumerate}
  \item \textbf{コマンドは自分の手で打つ。}
    コピー＆ペーストでも動きますが、手で打つと「どこで区切るのか」「何を指定しているのか」が体に入ります。
    長いもの（URL や鍵）だけはコピーして構いません。
  \item \textbf{分からない単語は、付録の用語集を引く。}
    それでも分からなければ、Claude に聞いてください（第6章で聞き方を説明します）。
  \item \textbf{30分詰まったら、三角に聞く。}
    30分考えて分からないことは、ひとりで3時間考えてもたいてい分かりません。
    LINE で聞くか、スプシの詰まりメモ欄に書いてください。聞き方の型は次の節で説明します。
\end{enumerate}

\section{進み具合の管理：スプレッドシート}

やることは、すべて三角が用意したGoogleスプレッドシート（以下「スプシ」）に1行ずつ並んでいます。
この教材のところどころに \sheet{1-3} のような印があり、スプシの行と対応しています。

\begin{center}\small
\begin{tabularx}{\linewidth}{>{\sffamily\bfseries}l X}
\toprule
列 & 書くこと \\
\midrule
ステータス & 未着手 → 進行中 → レビュー待ち → 完了。手をつけたら「進行中」にする \\
着手日・完了日 & その日の日付（\key{Ctrl}+\key{;} で今日の日付が入る） \\
PRリンク & PR を出したら URL を貼る（第2巻から） \\
詰まりメモ & 詰まったことを下の「質問の型」で書く。解決したら解決方法も追記する \\
三角確認 & \textbf{三角が付けるので触らない} \\
\bottomrule
\end{tabularx}
\end{center}

シートの一番上には、フェーズごとの完了率と「必須の5項目」が表示されます。
必須の5項目は、この教材全体を通して最終的にできるようになってほしいことです。

\section{質問の型}

三角に質問するときも、Claude に聞くときも、次の4つをそろえます。

\begin{tcolorbox}[kwbase,colback=washi,colframe=kon!40,title={質問の型（4点セット）},
  colbacktitle=kon!40]
\begin{description}[font=\sffamily\bfseries,leftmargin=6em,style=nextline]
  \item[やりたいこと] 何をしようとしていたか（例：STEP 3 で \code{git push} しようとした）
  \item[期待したこと] どうなるはずだったか（例：GitHub にファイルが表示される）
  \item[実際] 何が起きたか。\textbf{エラーメッセージは全文コピーして貼る}
  \item[試したこと] 自分で試したこと、調べたこと、Claude の回答
\end{description}
\end{tcolorbox}

「動きません」だけでは、三角にも Claude にも何が起きているのか分かりません。
4点セットを書いている途中で自分で答えに気づくことも、意外とよくあります。

\section{各フェーズの終わりにすること}

各フェーズの最後には、次の2つがあります。

\begin{description}[font=\sffamily\bfseries]
  \item[腕だめし] 手本どおりではなく、少しだけ違う課題です。
    \textbf{まず自分で「どのファイルの、どこを、どう変えれば、どうなるはずか」をメモしてから}、
    手を動かします（Claude を使ってもかまいません）。
  \item[説明タスク] 作ったものの仕組みを、三角に自分の言葉で説明します。
    定例の時間を使います。うまく説明できなかったところは宿題にして、次の定例で説明し直します。
\end{description}

\section{AI と一緒に作るということ}

この教材は、Claude Code（AI）に手伝ってもらいながら進める前提で書いてあります。
AI に書かせること自体は問題ありません。ただし、次の線だけは守ってください。

\begin{caution}[完了と言ってよい条件]
\begin{enumerate}
  \item \textbf{何を変えたか、なぜそう変えたかを、自分の言葉で説明できる}
  \item \textbf{動作を自分の目で確かめた}（どう確かめたかを書ける）
  \item \textbf{差分（変更前と変更後の違い）を自分で読んだ}
\end{enumerate}
この3つがそろわない状態の「動いたからOK」は、完了ではありません。
本番のサイトを保守するようになったとき、説明できないコードは直せないからです。
\end{caution}

\section{予定の目安}

期限はありません。週に5〜8時間ほど使えると、第1巻は1か月前後で終わる分量です。
定例は週1回、15〜30分です。進み具合の確認と、説明タスクに使います。
